\section{Landauer Theory}\label{ch:theory:sec:landauer_theory}
In order to make predictions about the electronic and thermoelectric properties of QDs, a good model for electronic transport in mesoscopic to molecular-scale devices is necessary. In the limit of predominantly SET and assuming coherent transport, extending the Landauer formalism (originally a model for electrical conduction) towards thermoelectric effects has proven successful and is widely used. The Landauer formalism is a concept formulated in 1947 by Rolf Landauer that models conduction as transmission. The basic assumption is that two electron reservoirs are connected by a scatterer (QD island for the transport measurements presented here) through which electron transport is possible. The transport is then determined by the probability of electrons moving from one reservoir to the other
through the scatterer. This probability is always between 0 and 1, which means there
is an upper limit on conductance.

Employing the Landauer approach, the current through such a structure can be
written as:
\begin{equation}
    I = \frac{2e}{h} \int\limits_{-\infty}^{+\infty}\mathcal{T}(\epsilon)[f_\textrm{L}(E) - f_\textrm{R}(E)]\text{d}\epsilon,
\end{equation}
where $\mathcal{T}$ is the transmission probability of an electron with single-particle energy $\epsilon$ to move across the source and drain leads. The Fermi-Dirac statistics
$f_\textrm{FD}(\epsilon, T) = \frac{1}{1+e^{\epsilon_i - \mu/k_BT}} $ details the average number of fermions (electrons in our case) in a single-particle state $i$ with energy $\epsilon_i$ in thermodynamic equilibrium at temperature $T$. Therefore, it is used to quantify how many electrons from the source leads have a certain energy to be able to cross the molecular junction: this event, as everything in quantum mechanics, is detailed by the amplitude of the transmission probability $\mathcal{T}(\epsilon)$.

One can calculate the number of electrons in a crystal by multiplying the density of
states $N(\epsilon)$ with the FD distribution and summing up this product for all energies:
\begin{equation}
    n = \int\limits_0^{\infty}f_\textrm{FD}(\epsilon)N(\epsilon)\text{d}\epsilon,
\end{equation}

The density of states $N(\epsilon)$ denotes the number of electron states per unit of energy and volume and differs for different materials.
In such a system the chemical potential is defined as the amount of energy needed
to add an additional electron to the system and naturally this quantity has an
inherent temperature dependence as the energy of the system changes with $T$. An
approximation for the chemical potential as a function of temperature and energy
has been developed using the Sommerfeld expansion (cf. Appendix~\ref{appendix:sommerfeld}): since we are interested in low-temperature properties, we will consider only terms up to the second order of Eq.~\ref{app-1:eq:coefficients_sommerfeld}:
\begin{equation}
    \mu \simeq \epsilon_{F}-\left.\frac{\pi^{2}}{6}\left(k_{B} T\right)^{2} \frac{1}{N\left(\epsilon_{F}\right)} \frac{\partial N(\epsilon)}{\partial \epsilon}\right|_{\epsilon=\epsilon_{F}}
\end{equation}
As is evident from this equation, the chemical potential decreases as the temperature
rises, resulting in a potential gradient whenever a temperature gradient is applied.
Any difference in $T$ will therefore lead to a flow of electrons as the cost of adding
another electron to the hot side would be lower than adding it to the cold side. This
results in the build-up of a thermoelectric voltage in an open circuit.

By multiplying this probability with the elementary charge $e$ gives the current for one energy level with a factor of two added due to spin degeneracy. After adding all of these possible transport channels together, which is done by taking the integral over all energies, and accounting for
dimension normalization via the Planck constant $h$ it is possible to calculate the total current $I$ across the junction.

The electrical conductance is then retrieved by taking the derivative of the current:

\begin{align}
    G & = -\frac{2e}{h} \int\limits_{-\infty}^{+\infty}\mathcal{T}(\epsilon)\pdv{f(T, \epsilon)}{\epsilon}\text{d}\epsilon,\nonumber \\
      & = - G_0 \int\limits_{-\infty}^{+\infty}\mathcal{T}(\epsilon)\pdv{f(T, \epsilon)}{\epsilon}\text{d}\epsilon,
\end{align}

Several computational scheme may be used to compute the set of transmission eigenvalues for a given molecular hamiltonian, ranging from simple tight-binding to non-equilibrium Green's functions. The interested reader is referred to the specific literature~\cite{cuevas2010molecular}.